% Раздел про программную реализацию.
% Здесь удобно показывать архитектуру и примеры конфигурации/кода.
Практическая реализация адаптивной системы построена как модульный программный
контур, где каждый компонент отвечает за отдельный этап образовательного цикла:
получение событий, аналитика профиля студента, выбор следующего задания,
генерация объяснения и публикация результата в пользовательский интерфейс.
Такой подход позволяет развивать систему итеративно, не нарушая стабильность
всей платформы. В инженерных терминах архитектура организована вокруг
контрактов обмена сообщениями \texttt{event\_payload}, \texttt{state\_snapshot}
и \texttt{recommendation\_packet}, что делает взаимодействие между сервисами
прозрачным и тестируемым.

Функционально в системе выделяются четыре основные подсистемы. Первая
подсистема -- контур сбора данных, принимающий события от LMS и интерфейсов
пользователей. Вторая подсистема -- аналитический слой, в котором выполняется
очистка данных, агрегация признаков и обновление модели знания. Третья
подсистема -- рекомендательный модуль, принимающий решение о следующем шаге
обучения. Четвертая подсистема -- интерфейсный слой, обеспечивающий выдачу
рекомендации, отображение объяснения и сбор обратной связи от студента и
эксперта. Дополнительно используется слой хранения, где разделяются
операционные данные и аналитические срезы для отчетности.

Для повышения воспроизводимости развертывания параметры системы вынесены в
конфигурационный профиль курса. Это позволяет запускать одинаковый программный
контур в разных дисциплинах, изменяя только предметные настройки. Ключевыми
параметрами являются \texttt{mastery\_threshold}, \texttt{epsilon\_start},
\texttt{epsilon\_end}, \texttt{reward\_weights} и \texttt{max\_review\_gap}.
Значения подбираются на основании пилотной валидации и могут корректироваться
методистом при накоплении новых данных.

Архитектура демонстрационного контура показана на рисунке
\ref{fig:system-architecture}. Диаграмма отражает как основной поток данных,
так и контур обратной связи, в котором преподаватель и экспертные интерфейсы
влияют на дальнейшую настройку политики.

Для прозрачной настройки курса ключевые параметры адаптивной политики собраны в
таблице \ref{tab:policy-params}. Такой формат упрощает нормоконтроль и
согласование параметров между методистом и разработчиком.

\begin{table}[htbp]
  \caption{Ключевые параметры адаптивной политики}
  \label{tab:policy-params}
  \centering

  \begin{tabular}{|p{0.27\textwidth}|p{0.20\textwidth}|p{0.40\textwidth}|}
    \hline
    Параметр & Типовое значение & Назначение \\
    \hline
    \texttt{mastery\_threshold} & 0.78 & Порог перехода к более сложному шагу \\
    \hline
    \texttt{epsilon\_start} & 1.00 & Стартовый уровень исследовательского поведения политики \\
    \hline
    \texttt{epsilon\_end} & 0.05 & Нижняя граница доли исследовательских действий \\
    \hline
    \texttt{reward\_weights} & [0.45, 0.25, 0.20, 0.10] & Баланс учебного прогресса, вовлеченности, утомления и фрустрации \\
    \hline
    \texttt{max\_review\_gap} & 5 шагов & Максимальный интервал до повторной проверки навыка \\
    \hline
  \end{tabular}
\end{table}

\begin{figure}[htbp]
  \centering

  \begin{tikzpicture}[
    x=11mm,
    y=10mm,
    >={Stealth[length=2.4mm]},
    module/.style={draw, rounded corners, minimum width=22mm, minimum height=11mm, text width=21mm, align=center, fill=teal!8},
    data/.style={draw, trapezium, trapezium left angle=70, trapezium right angle=110, minimum width=20mm, minimum height=9mm, text width=19mm, align=center, fill=orange!15},
    line/.style={->, thick}
    ]
    \node[data] (events) at (0, 0) {Учебные\\события};
    \node[module] (collector) at (2.8, 0) {Слой\\сбора};
    \node[module] (analytics) at (5.6, 0) {Аналитика\\состояния};
    \node[module] (recsys) at (8.4, 0) {Рекомендательный\\модуль};
    \node[module] (ui) at (11.2, 0) {Интерфейс\\преподавателя\\и студента};
    \node[data] (profile) at (5.6, -2.6) {Профиль\\обучающегося};

    \draw[line] (events) -- (collector);
    \draw[line] (collector) -- (analytics);
    \draw[line] (analytics) -- (recsys);
    \draw[line] (recsys) -- (ui);
    \draw[line] (analytics) -- (profile);
    \draw[line] (profile) -- (recsys);
    \draw[line, dashed] (ui.south west) to[out=235,in=-20,looseness=1.25] node[pos=0.52, below, yshift=-2pt, draw=none, fill=white, inner sep=1.2pt, font=\footnotesize] {экспертная обратная связь} (collector.south east);
  \end{tikzpicture}

  \caption{Модульная архитектура адаптивной системы обучения}
  \label{fig:system-architecture}
\end{figure}

В процессе разработки использовался принцип "контракт сначала": каждый сервис
публикует явное описание входов и выходов. Это особенно важно для адаптивных
систем, где логика принятия решения должна быть аудитируемой. Например, выход
рекомендательного модуля включает не только идентификатор задания, но и
диагностические поля \texttt{policy\_score}, \texttt{confidence} и
\texttt{reason\_tags}. Благодаря этому в журнале можно восстановить цепочку
принятия решения и сопоставить ее с фактическим результатом выполнения задания.

Ниже приведен пример конфигурации рекомендательного модуля, используемой для
запуска базового сценария.

\begin{lstlisting}[caption={Фрагмент конфигурации рекомендательного модуля}]
policy = {
    "algorithm": "dqn",
    "epsilon_start": 1.0,
    "epsilon_end": 0.05,
    "gamma": 0.95,
    "batch_size": 64,
    "target_sync_steps": 200,
}
\end{lstlisting}

Для устойчивой обработки событий применяется этап нормализации входных данных.
На этом этапе удаляются технические дубликаты, проверяются обязательные поля и
рассчитываются производные признаки, включая латентность ответа и число
повторных попыток.

\begin{lstlisting}[style=sthgcode-python,caption={Пример нормализации входного учебного события}]
def normalize_event(raw_event: dict) -> dict:
    required = {"student_id", "task_id", "is_correct", "started_at", "finished_at"}
    missing = required - set(raw_event)
    if missing:
        raise ValueError(f"Missing fields: {sorted(missing)}")

    latency = max(0, raw_event["finished_at"] - raw_event["started_at"])
    return {
        "student_id": int(raw_event["student_id"]),
        "task_id": int(raw_event["task_id"]),
        "is_correct": bool(raw_event["is_correct"]),
        "latency_sec": float(latency),
        "attempt_no": int(raw_event.get("attempt_no", 1)),
    }
\end{lstlisting}

Инфраструктурный контур реализует отдельный API для экспертной обратной связи.
Такой API нужен, чтобы преподаватель мог корректировать оценку рекомендации и
уточнять текст объяснения без ручного редактирования базы данных.

\begin{lstlisting}[style=sthgcode-java,caption={Пример обработчика экспертной корректировки рекомендации}]
@PostMapping("/api/expert/review")
public ReviewResult submitExpertReview(@RequestBody ExpertReview payload) {
    recommendationService.applyRewardOverride(
        payload.recommendationId(),
        payload.deltaReward(),
        payload.comment()
    );
    explanationService.updateExplanation(
        payload.recommendationId(),
        payload.editedExplanation()
    );
    return ReviewResult.ok(payload.recommendationId());
}
\end{lstlisting}

Распределение ролей и типовых пользовательских сценариев приведено в таблице
\ref{tab:user-roles-scenarios}. Таблица фиксирует, какие сущности данных
задействуются в каждом контуре взаимодействия.

\begin{table}[htbp]
  \caption{Роли пользователей и ключевые сценарии взаимодействия}
  \label{tab:user-roles-scenarios}
  \centering

  \begin{tabular}{|p{0.18\textwidth}|p{0.34\textwidth}|p{0.35\textwidth}|}
    \hline
    Роль & Основной сценарий & Ключевые поля/сущности \\
    \hline
    Студент & Получение следующего задания и объяснения выбора & \texttt{task\_id}, \texttt{confidence}, \texttt{reason\_text} \\
    \hline
    Методист & Публикация и ревизия контента курса & \texttt{skill\_map}, \texttt{difficulty\_band}, \texttt{prerequisite\_links} \\
    \hline
    Эксперт & Коррекция подкрепления и редактирование объяснений & \texttt{recommendation\_id}, \texttt{expert\_delta\_reward}, \texttt{edited\_reason} \\
    \hline
    Администратор & Мониторинг стабильности и запуск регламентных процедур & \texttt{policy\_version}, \texttt{drift\_alerts}, \texttt{deployment\_status} \\
    \hline
  \end{tabular}
\end{table}

Отдельная задача проектирования -- обеспечить предсказуемую эксплуатацию модели
в условиях частично наблюдаемых данных. В реальных курсах часть событий может
приходить с задержкой, а часть метрик пересчитывается асинхронно. Чтобы
избежать нестабильного поведения, в системе используется буфер обновления
состояния: решения принимаются только по согласованному срезу данных, а
задержанные события попадают в очередь компенсационного пересчета. На уровне
конфигурации это задается флагами \texttt{stale\_event\_policy} и
\texttt{recompute\_window}.

Для поддержки CI/CD предусмотрен минимальный технический скрипт, который
подготавливает окружение и запускает обязательные проверки перед выкладкой
новой версии политики.

\begin{lstlisting}[style=sthgcode-python,caption={Пример автоматизированного pre-deploy сценария проверки}]
import subprocess


def run(cmd: str) -> None:
    subprocess.run(cmd, shell=True, check=True)


run("make check-style")
run("make build-matrix")
run("make check-fonts")
run("make check-structure")
\end{lstlisting}

При переносе прототипа в учебную эксплуатацию особое внимание уделяется
вопросам данных и безопасности. Для персональных образовательных профилей
используется разделение идентифицирующих и аналитических данных: в
аналитическом контуре студент представлен техническим идентификатором, а
операции с персональными сведениями ограничены административным слоем с
журналированием доступа. Такой подход уменьшает риск неконтролируемого
использования чувствительной информации и одновременно сохраняет возможность
корректной статистической обработки учебного поведения. На уровне сервисов это
обеспечивается политиками \texttt{data\_retention\_days},
\texttt{anonymize\_exports} и \texttt{pii\_scope}.

Отдельный контур реализован для мониторинга качества рекомендаций после
обновления модели. В практике внедрения адаптивных систем важно не только
получить хороший результат в офлайн-оценке, но и контролировать стабильность
поведения в реальном потоке студентов. Для этого в системе фиксируются метрики
дрейфа: изменение распределения признаков, отклонение средней латентности
ответа, рост доли ручных экспертных коррекций и изменение структуры отказов
студентов от рекомендованных шагов. Если один из индикаторов превышает заданный
порог, политика переводится в режим осторожной эксплуатации с пониженным
\texttt{exploration\_rate}, а ответственному методисту отправляется уведомление
для ревизии контента и параметров.

Кроме технического мониторинга, предусмотрен регламент предметной валидации.
Перед публикацией новой версии курса проводится совместная проверка методиста и
преподавателя: анализируются предпосылки тем, корректность карты навыков,
соответствие сложности задач заявленным результатам обучения и качество
автоматических объяснений. По итогам ревизии обновляются шаблоны обратной связи
и матрица весов награды. Такая процедура снижает риск того, что система будет
оптимизировать формальные метрики в ущерб дидактической целесообразности. В
результате программный контур сохраняет управляемость не только на уровне кода,
но и на уровне образовательной методики.

В совокупности описанные решения формируют инженерную основу адаптивной
платформы, где алгоритмическая логика встроена в управляемый и проверяемый
программный процесс. Такой контур облегчает масштабирование и снижает риск
"скрытых" регрессий, которые в образовательной системе могут приводить не
только к техническим сбоям, но и к методическим ошибкам в выборе траектории
обучения.
